\documentclass[aspectratio=169,10pt,spanish]{beamer}

\input{../../../_preambulo_beamer.tex}
\input{../../../_comandos_beamer.tex}

\selectlanguage{spanish}

\title{MA1001B - Modelación Estadística para la Toma de Decisiones}
\subtitle{Lección 01: Población, muestras y tipos de datos}
\author{}
\institute{Tecnol\'ogico de Monterrey}
\date{}

\begin{document}

\begin{frame}[plain]
  \vspace{1.2cm}
  {\Large\bfseries MA1001B - Modelación Estadística para la Toma de Decisiones\par}
  \vspace{0.7cm}
  {\large Lección 01: Población, muestras y tipos de datos\par}
  \vspace{0.7cm}
  {\color{TecRojo}\rule{\textwidth}{1.2pt}}
  \vspace{0.8cm}

  Facultad MA1001B\\[0.3cm]
  Periodo\\[0.3cm]
  Tecnol\'ogico de Monterrey
\end{frame}

\begin{frame}{La pregunta estadística}
  \begin{alertblock}{Pregunta guía}
    ¿Cómo puede un conjunto limitado de observaciones sustentar una afirmación
    sobre una población de negocio más grande?
  \end{alertblock}

  \vspace{0.6em}
  Al final de esta lección, usted puede:
  \begin{itemize}
    \item definir la población objetivo y la muestra observada;
    \item conectar un parámetro poblacional con su estadístico muestral;
    \item clasificar variables según su significado analítico;
    \item identificar por qué el tamaño de muestra no elimina el sesgo de selección.
  \end{itemize}

  \vfill
  \scriptsize Todos los registros usados hoy son sintéticos. No se usan datos reales de empresa.
\end{frame}

\begin{frame}{Población, muestra, parámetro y estadístico}
  \small
  \begin{center}
    \begin{tabular}{p{0.19\textwidth}p{0.32\textwidth}p{0.32\textwidth}}
      \toprule
      \textbf{Rol} & \textbf{Población} & \textbf{Muestra} \\
      \midrule
      Observaciones & Todas las unidades definidas por la pregunta & Unidades realmente seleccionadas \\
      Tamaño & $N$ & $n$ \\
      Resumen numérico & Parámetro, como $\mu$ & Estadístico, como $\bar{x}$ \\
      Estatus & Fijo para la población definida & Cambia entre muestras \\
      \bottomrule
    \end{tabular}
  \end{center}

  \vspace{0.6em}
  \begin{exampleblock}{Pregunta sintética de operaciones}
    ¿Cuál es el tiempo medio de procesamiento de los $12{,}000$ pedidos del
    periodo? Una muestra aleatoria de 200 pedidos proporciona $\bar{x}$ como
    estimación de la media poblacional $\mu$.
  \end{exampleblock}
\end{frame}

\begin{frame}[fragile]{Paso 1: Una muestra estima una cantidad poblacional}
  \begin{columns}[T]
    \begin{column}{0.50\textwidth}
      \small
      \begin{lstlisting}
population = build_population()
sample = draw_random_sample(population)

mu = population[variable].mean()
x_bar = sample[variable].mean()
error = abs(x_bar - mu)
      \end{lstlisting}

      \begin{block}{Salida verificada}
        $N=12{,}000$ y $n=200$\\
        $\mu=45.26$ minutos\\
        $\bar{x}=46.11$ minutos\\
        Error absoluto $=0.85$ minutos
      \end{block}
    \end{column}
    \begin{column}{0.50\textwidth}
      \centering
      \includegraphics[width=\textwidth]{../figuras/fundos_datos_01_poblacion_muestra.png}
      \vspace{0.2em}
      \scriptsize Distribuciones de población y muestra aleatoria del Paso 1
    \end{column}
  \end{columns}
\end{frame}

\begin{frame}{Tipos de datos de negocio}
  \small
  \begin{center}
    \begin{tabular}{p{0.23\textwidth}p{0.28\textwidth}p{0.34\textwidth}}
      \toprule
      \textbf{Tipo} & \textbf{Pregunta} & \textbf{Ejemplo} \\
      \midrule
      Categórica & ¿Qué grupo o etiqueta? & Canal, región \\
      Cuantitativa discreta & ¿Cuántas unidades contables? & Artículos por pedido \\
      \shortstack[l]{Cuantitativa\\continua} & ¿Cuánto se midió? & Tiempo, valor del pedido \\
      \bottomrule
    \end{tabular}
  \end{center}

  \vspace{0.7em}
  \begin{alertblock}{El significado analítico controla la clasificación}
    \texttt{order\_id} contiene dígitos, pero la aritmética sobre un
    identificador no tiene sentido de negocio. Actúa como etiqueta.
  \end{alertblock}
\end{frame}

\begin{frame}[fragile]{Paso 2: Una tabla contiene varios tipos de datos}
  \begin{columns}[T]
    \begin{column}{0.42\textwidth}
      \small
      \begin{lstlisting}
roles = {
  "channel": "Categórica",
  "items_per_order":
      "Cuantitativa discreta",
  "processing_time_minutes":
      "Cuantitativa continua",
}
      \end{lstlisting}

      La misma muestra de 200 filas admite resúmenes distintos porque cada
      variable responde a un tipo distinto de pregunta.
    \end{column}
    \begin{column}{0.58\textwidth}
      \centering
      \includegraphics[width=\textwidth]{../figuras/fundos_datos_02_tipos_datos.png}
      \vspace{0.2em}
      \scriptsize Conteos para categorías y discretas; histograma para un
      valor continuo medido.
    \end{column}
  \end{columns}
\end{frame}

\begin{frame}{Paso 3: Una muestra sesgada más grande tiene más error}
  \begin{columns}[T]
    \begin{column}{0.43\textwidth}
      \small
      Parámetro poblacional: $\mu=45.26$ minutos

      \vspace{0.5em}
      \begin{tabular}{lrr}
        \toprule
        \textbf{Muestra} & \textbf{$n$} & \textbf{Error} \\
        \midrule
        Aleatoria & 200 & 0.85 min \\
        Conveniencia & 1{,}000 & 6.24 min \\
        \bottomrule
      \end{tabular}

      \vspace{0.7em}
      La muestra de conveniencia contiene solo pedidos En línea. Su mayor
      tamaño no restaura los canales faltantes.

      \vspace{0.5em}
      \textbf{Límite:} este ejemplo determina un sesgo de selección posible.
      No afirma que toda muestra de conveniencia tenga este error exacto.
    \end{column}
    \begin{column}{0.57\textwidth}
      \centering
      \includegraphics[width=\textwidth]{../figuras/fundos_datos_03_sesgo_muestreo.png}
      \vspace{0.2em}
      \scriptsize Tamaño de muestra y error de estimación de la misma población
    \end{column}
  \end{columns}
\end{frame}

\begin{frame}{Verificación de conceptos}
  \begin{enumerate}
    \item Un analista estudia todos los pedidos del mes e inspecciona 300.
      Identifique la población y la muestra.
    \item Los 300 pedidos tienen un valor medio de MXN 1{,}840. ¿Es un
      parámetro o un estadístico? Explique.
    \item Clasifique segmento de cliente, número de artículos y tiempo de entrega.
    \item ¿Por qué 1{,}000 registros convenientemente disponibles pueden dar
      una peor estimación que 200 seleccionados al azar?
  \end{enumerate}

  \vspace{0.6em}
  \begin{block}{Conexión con la Lección 02}
    La sesión puede continuar directamente con muestreo aleatorio simple,
    estratificado y por conglomerados.
  \end{block}

  \vfill
  \scriptsize Fuente: OpenStax, \textit{Introductory Business Statistics 2e},
  Capítulo 1, Secciones 1.1 y 1.2. CC BY-NC-SA.
\end{frame}

\appendix



\end{document}
